\documentclass[11pt,a4paper]{article}

% Paquetes Requeridos 
\usepackage[utf8]{inputenc}
\usepackage[spanish,es-tabla]{babel}
\usepackage[margin=2.5cm]{geometry}
\usepackage{amsmath}
\usepackage{amssymb}
\usepackage{amsfonts}
\usepackage{booktabs}
\usepackage{graphicx}
% Usar estilo APA y compatibilidad con \bibliographystyle{apacite}
\usepackage[round,authoryear]{natbib}
\usepackage{microtype}
% Mejora de índices y control de reejecuciones de compilación
\usepackage{url}
\usepackage{array}
\usepackage{hyperref}
\usepackage{bookmark}

% Configuración de Hipervínculos
\hypersetup{
    colorlinks=true,
    linkcolor=blue,
    citecolor=blue,
    urlcolor=blue
}

% Título y Datos del Autor
\title{\textbf{Análisis de datos de panel del balance hídrico en Paraguay: Cuantificación de la sequía y la aridificación ante el calentamiento global (2000--2023)}}
\author{
    \textbf{Hernán J. Vargas Ojeda}\thanks{Departamento de Estadística, Facultad de Ciencias Exactas y Naturales, Universidad Nacional de Asunción, San Lorenzo, Paraguay. Correo electrónico: \href{mailto:hjvargaso@gmail.com}{hjvargaso@gmail.com}} \\
    \small Orientador: Prof. Dr. Gonzalo Martin Vicente* \\
    \small Co-Orientador: Prof. Dr. Gustavo Ignacio Rivas M.*
}
\date{\small Mayo de 2026}

\begin{document}


\maketitle

\begin{abstract}
\noindent Este artículo evalúa la evolución espacio-temporal de las sequías y los procesos de aridificación en Paraguay durante el período 2000--2023, analizando el impacto de la demanda evaporativa en un contexto de calentamiento global. Se estructuró una base de datos en formato de panel balanceado a partir de los registros diarios de la red oficial de monitoreo terrestre. Se implementó un algoritmo de imputación por capas térmicas jerárquicas y un criterio de no imputación para la precipitación. El territorio nacional fue regionalizado mediante técnicas de agrupamiento jerárquico aglomerativo (HAC) y optimizado mediante el coeficiente de silueta promedio, identificando dos subregiones climáticas homogéneas de balance hídrico: la subregión sureste (Clúster 1) y la subregión centrochaco (Clúster 2). Las estimaciones de tendencias estacionales y regionales se realizaron mediante un modelo de regresión de datos de panel de efectos aleatorios con errores estándar de Driscoll y Kraay (1998) robustos a la correlación serial y a la fuerte dependencia espacial confirmada por la prueba de Pesaran. Los resultados revelan una marcada divergencia estacional en el balance hídrico: mientras las estaciones de otoño e invierno se mantuvieron hídricamente estables, las épocas cálidas de primavera y verano registraron trayectorias de aridificación. La tasa de aridificación estival más intensa del país se concentró en la subregión agropecuaria del sureste (Clúster 1). El contraste paralelo entre el modelo de balance climático (SPEI-6) y el de precipitación univariada (SPI-6) demostró la validez del rol del estrés térmico; debido a que el régimen de lluvias permaneció estadísticamente estable (SPI-6 no significativo), la significancia de la diferencia de pendientes estimada en el verano del Clúster 2 al 5\% y en el verano del Clúster 1 al 10\% son evidencias de que las sequías en Paraguay están adquiriendo una naturaleza predominantemente térmica en el siglo XXI, donde la evaporación forzada por el incremento sistemático de las temperaturas máximas extremas induce procesos de aridificación estacional crónicos incluso bajo volúmenes de precipitación estables.

\noindent \textbf{Palabras clave:} Datos de panel, balance hídrico, aridificación estacional, SPEI, efectos aleatorios, Driscoll-Kraay, calentamiento global, Paraguay.
\end{abstract}

\newpage
\renewcommand{\abstractname}{Abstract}
\begin{abstract}
\noindent This article evaluates the spatio-temporal evolution of droughts and desertification processes in Paraguay during the period 2000-2023, analyzing the impact of evaporative demand in a context of global warming. A balanced panel database was structured from daily records of the official terrestrial monitoring network. A hierarchical thermal layer imputation algorithm and a non-imputation criterion for precipitation were implemented. The national territory was regionalized using agglomerative hierarchical clustering (HAC) techniques and optimized using the average silhouette coefficient, identifying two homogeneous climatic subregions of water balance: the southeastern subregion (Cluster 1) and the central Chaco subregion (Cluster 2). Seasonal and regional trend estimates were performed using a random-effects panel data regression model with Driscoll and Kraay (1998) robust standard errors to serial correlation and strong spatial dependence confirmed by Pesaran's test. The results reveal a marked seasonal divergence in the water balance: while autumn and winter seasons remained hydrologically stable, the warm spring and summer periods recorded desertification trajectories. The most intense summer desertification rate in the country was concentrated in the agricultural subregion of the southeast (Cluster 1). The parallel contrast between the climate balance model (SPEI-6) and the univariate precipitation model (SPI-6) demonstrated the validity of the role of thermal stress; because the rainfall regime remained statistically stable (SPI-6 not significant), the significance of the estimated slope difference in summer of Cluster 2 at 5\% and in summer of Cluster 1 at 10\% are evidence that droughts in Paraguay are acquiring a predominantly thermal nature in the 21st century, where evaporation forced by the systematic increase of extreme maximum temperatures induces chronic seasonal desertification processes even under stable precipitation volumes.

\noindent \textbf{Key words:} Panel data, water balance, seasonal aridification, SPEI, random effects, Driscoll-Kraay, global warming, Paraguay.
\end{abstract}

\section{Introducción}

En todo balance hídrico, los ingresos están representados por las precipitaciones y los egresos por la evapotranspiración potencial (ETP), entendida como la cantidad de agua que la atmósfera intenta extraer del suelo y la vegetación. Tradicionalmente, la sequía se ha conceptualizado como un fenómeno impulsado exclusivamente por la escasez de lluvias. No obstante, el calentamiento global inequívoco del sistema climático \citep{stocker_climate_2013} ha redefinido este paradigma, demostrando que el incremento sostenido de las temperaturas exacerba la demanda evaporativa de la atmósfera, induciendo déficits hídricos crónicos aún bajo regímenes pluviométricos estables, un fenómeno denominado por la literatura moderna como sequías más cálidas(hotter droughts) \citep{trenberth_global_2014}.

A nivel regional, la Cuenca del Río de la Plata experimentó entre 2019 y 2022 una de las sequías plurianuales más severas de su registro instrumental, provocando estiajes históricos en los ríos Paraguay y Paraná con pérdidas millonarias en el sector agropecuario, la navegación comercial y la generación hidroeléctrica cite{naumann2023}. A pesar de que investigaciones nacionales han documentado tendencias ascendentes y estadísticamente significativas en las temperaturas medias de Paraguay, aún no se ha abordado de manera sistemática y cuantitativa la sequía y la aridificación en un contexto de cambio climático global.

El problema fundamental radica en que este vacío de conocimiento genera una incertidumbre climática que se traduce en riesgos no cuantificados para la seguridad hídrica y la estabilidad económica del país. Esta investigación aborda dicho vacío bajo dos dimensiones analíticas: por un lado, examinando el vínculo directo de la temperatura con el aumento de la demanda evaporativa; por otro, analizando su comportamiento estacional y geográfico mediante un modelado de datos de panel robusto que considera la fuerte correlación espacial sincrónica de las estaciones meteorológicas del país.

El objetivo del trabajo es evaluar la evolución del balance hídrico en Paraguay mediante el análisis integrado de precipitación y temperatura, con el fin de caracterizar los eventos de sequía, identificar patrones espaciales y estacionales, y detectar posibles dinámicas de aridificación a largo plazo. En términos metodológicos, el estudio construye una base climática en formato de datos de panel, regionaliza el territorio paraguayo a partir de la variabilidad común del balance hídrico y ajusta modelos de datos de panel para cuantificar tendencias temporales del SPEI según regiones y estaciones climáticas.

De este modo, la investigación contribuye a la literatura climática nacional en tres sentidos. Primero, ofrece una caracterización sistemática de la sequía en Paraguay incorporando una métrica sensible al efecto térmico. Segundo, permite distinguir entre déficits asociados únicamente a la precipitación y déficits derivados del balance entre precipitación y demanda evaporativa. Tercero, propone un marco metodológico replicable para el análisis de sequía y aridificación en contextos territoriales donde la disponibilidad de series climáticas permite estructurar datos
longitudinales y espaciales comparables.

\section{Marco conceptual y enfoque analítico}

\subsection{Sequía, aridificación y calentamiento global}

La sequía es un fenómeno climático de desarrollo lento, acumulativo y espacialmente heterogéneo, caracterizado por una disponibilidad de agua inferior a las condiciones normales de un territorio durante un período determinado. Su análisis requiere considerar no solo la intensidad del déficit hídrico, sino también su duración, frecuencia y extensión espacial \citep{mishra2010,wmo2017}. En términos generales, la literatura distingue entre sequía meteorológica, agrícola, hidrológica y socioeconómica, según el componente del sistema afectado y la escala temporal de manifestación del déficit \citep{mishra2010,velasco2005sequia}.

La sequía meteorológica se vincula principalmente con anomalías negativas de precipitación. Sin embargo, en contextos de calentamiento global, esta aproximación puede resultar incompleta, ya que el incremento de la temperatura aumenta la demanda evaporativa de la atmósfera y puede intensificar el déficit hídrico aun cuando la precipitación no presente reducción significativa \citep{trenberth_global_2014,massondelmotte2021}. Por esta razón, el análisis contemporáneo de la sequía requiere integrar simultáneamente la oferta hídrica, representada por la precipitación, y la demanda atmosférica de agua, aproximada mediante la evapotranspiración potencial.

La aridificación, a diferencia de la sequía, no se refiere necesariamente a un evento transitorio, sino a una trayectoria climática de largo plazo hacia condiciones más secas, cuando mediante un proceso dinámico de largo
plazo se produce un cambio estructural hacia un clima más árido entonces se considera una aridificación \citep{vicenteserrano2010}. La sequía puede actuar como un factor desencadenante o un acelerador crítico del proceso de aridificación. 

Estas problemáticas son relevantes para Paraguay, donde la ocurrencia de eventos secos recientes, las tendencias térmicas ascendentes y la alta dependencia del sector agropecuario configuran un escenario de elevada vulnerabilidad climática \citep{cepal2014}. Como afirma \citet{mishra2010}, estos fenómenos comprometen de manera
directa la seguridad alimentaria al afectar directamente a la agricultura, el acceso al recurso hídrico y la estabilidad macroeconómica.





\section{Materiales y Métodos}

\subsection{Fuentes de datos y variables}
Los datos primarios corresponden a los registros diarios de precipitación acumulada ($P_t$, mm), temperatura media diaria ($T_{\mathrm{med},t}$, $^\circ$C), temperatura máxima ($T_{\mathrm{max},t}$, $^\circ$C) y temperatura mínima ($T_{\mathrm{min},t}$, $^\circ$C) para el período comprendido entre el 1 de enero de 2000 y el 31 de diciembre de 2023, provistos por la Dirección de Meteorología e Hidrología (DMH) de Paraguay. La base original contenía 19 estaciones meteorológicas terrestres. Tras un análisis de completitud temporal, se determinó que la estación de Paraguarí (PARA) presentaba un truncamiento severo (pérdida masiva del 26\% de la serie), por lo que fue excluida definitivamente. Las 18 estaciones restantes exhibieron tasas de completitud de fechas superiores al 97.9\%, conformando la base del panel balanceado inicial de 157,788 registros diarios.

\subsection{Procesamiento, depuración e imputación de datos}
Para garantizar la continuidad de la serie de temperatura media diaria ($T_{\text{med}}$) y evitar la pérdida de observaciones mensuales, se diseñó e implementó un algoritmo de imputación estructurado en tres capas jerárquicas:
\begin{enumerate}
    \item \textbf{Capa 1: Imputación por relación térmica.} En días con omisión de $T_{\text{med}}$, pero con registros válidos de temperaturas extremas, se estimó el valor como el promedio aritmético: $T_{\text{med},t} = (T_{\text{mín},t} + T_{\text{máx},t})/2$ (recuperando 342 días).
    \item \textbf{Capa 2: Interpolación lineal temporal.} Para vacíos consecutivos cortos ($\le 3$ días), se aplicó una interpolación lineal temporal basada en la inercia térmica de los días adyacentes válidos (recuperando 655 días).
    \item \textbf{Capa 3: Imputación espacial por base vecina.} Para vacíos prolongados ($\ge 4$ días), se calculó una matriz de distancias ortodrómicas geográficas entre estaciones para identificar unívocamente la estación vecina más cercana y transferir directamente el registro térmico diario correspondiente (recuperando 1,114 observaciones). 
\end{enumerate}

Al término de la tercera capa, la tasa de completitud global en las series de temperatura diaria alcanzó el 99.99\%. En contraposición, para la variable precipitación se adoptó la decisión metodológica de no aplicar ningún método de imputación. Debido a la intermitencia y asimetría de la lluvia, forzar estimaciones espaciales o temporales introduciría sesgos artificiales en la frecuencia de días lluviosos y la magnitud acumulada de los eventos \citep{wooldridge2010}.

\subsection{Agregación mensual y cálculo de índices climatológicos}
El panel diario depurado de 18 estaciones se agregó a una frecuencia mensual de 288 meses por estación ($N \times T = 5,184$ observaciones mensuales potenciales), aplicando la regla de exclusión de la Organización Meteorológica Mundial \citep{wmo2017}: una variable agregada mensual se considera nula ($NaN$) si registra menos de 25 días con observaciones válidas. 

\subsubsection{Cálculo del SPEI-6}
El cálculo del Índice de Precipitación-Evapotranspiración Estandarizado a escala de 6 meses (SPEI-6) se ejecutó secuencialmente de la siguiente manera:
\begin{enumerate}
    \item \textit{Cálculo de la ETP mensual (Thornthwaite):} Estimada según el método empírico de \citet{thornthwaite_approach_1948}, adaptado al hemisferio sur mediante un factor de corrección de luz diurna ($K$) y sumando de forma acumulativa anual el índice de calor mensual ($i$) para la estabilidad del exponente $a$.
    \item \textit{Balance hídrico mensual ($D_i$):} Definido como $D_i = P_i - ETP_i$.
    \item \textit{Acumulación multiescalar ($D^k_i$):} El balance se acumuló para una ventana temporal de $k=6$ meses: $D^6_i = \sum_{j=0}^{5} D_{i-j}$.
    \item \textit{Ajuste paramétrico Pearson Tipo III:} Para modelar el balance hídrico acumulado de manera robusta ante la presencia de asimetrías y números negativos, los datos se ajustaron a una distribución teórica Pearson Tipo III mediante el método de máxima verosimilitud en Python (\textit{SciPy}).
    \item \textit{Estabilización numérica (Recorte de CDF):} Para evitar la generación de valores de probabilidad acumulada extremos ($0.0$ o $1.0$) en los extremos de la distribución, que se traducen en valores de anomalía infinitos ($\pm \infty$), se aplicó un recorte (\textit{clipping}) sobre la CDF calculado como:
    \begin{equation}
        P_{\text{ajustada}} = \max\left(\epsilon, \min\left(1 - \epsilon, F(x)\right)\right)
    \end{equation}
    Estableciendo un límite de tolerancia computacional de $\epsilon = 10^{-6}$.
    \item \textit{Estandarización Normal:} El valor del SPEI-6 final se obtuvo transformando la probabilidad $P_{\text{ajustada}}$ mediante la función cuantil normal estándar ($\Phi^{-1}$). El mismo procedimiento paramétrico, omitiendo la ETP y aplicando un ajuste de distribución Gamma para la precipitación acumulada, se utilizó para estimar la serie mensual del SPI-6.
\end{enumerate}

\subsection{Regionalización y modelado econométrico de datos de panel}
Para delimitar subregiones climáticas homogéneas de sequía sin imponer límites políticos arbitrarios, se aplicó un algoritmo de agrupamiento jerárquico aglomerativo (HAC) sobre las series de SPEI-6, utilizando la distancia de disimilitud euclídea y el criterio de enlace de varianza mínima de Ward. El número óptimo de clústeres ($k^*$) se determinó mediante la maximización del Coeficiente de Silueta Promedio de Rousseeuw \citep{rousseeuw1987}.

Una vez regionalizado el país, se especificó un modelo de regresión lineal para datos de panel con interacciones múltiples de la forma:
\begin{equation}
\begin{split}
    Y_{it} = & \beta_0 + \beta_1 T_t + \sum_{s=1}^{3} \gamma_s E_{s,t} + \delta_1 R_i \\
    & + \sum_{s=1}^{3} \theta_s (T_t \times E_{s,t}) + \phi_1 (T_t \times R_i) + \sum_{s=1}^{3} \psi_s (E_{s,t} \times R_i) \\
    & + \sum_{s=1}^{3} \omega_s (T_t \times E_{s,t} \times R_i) + u_{it}
\end{split}
\label{eq:modelo_completo}
\end{equation}

Donde $Y_{it}$ es el SPEI-6 (o SPI-6) de la estación $i$ en el mes $t$; $T_t$ es la tendencia mensual continua ($t = 1, \dots, 288$); $E_{s,t}$ son variables dicotómicas estacionales (otoño, primavera, verano, tomando el invierno como referencia); $R_i$ es la variable dicotómica regional y $u_{it} = \alpha_i + \varepsilon_{it}$ es el término de error compuesto. 

La elección entre efectos fijos (FE) y efectos aleatorios (RE) se evaluó mediante la prueba de Hausman (1978). La dependencia de sección cruzada (correlación espacial de los residuos provocada por choques macroclimáticos de gran escala como el ENSO) se diagnosticó formalmente mediante la prueba de Pesaran (2004). Para garantizar la validez de la inferencia ante autocorrelación serial temporal y dependencia espacial, se prescindió de las matrices de covarianza clásicas y se estimó el modelo final mediante la regresión de Driscoll y Kraay (1998).

Finalmente, las tendencias netas mensuales se calcularon mediante combinaciones lineales de los coeficientes del modelo ($\beta_{\text{neta},s} = \mathbf{w}_s' \boldsymbol{\hat{\beta}}$), validando su significancia mediante pruebas de Wald robustas utilizando la matriz de covarianza de Driscoll-Kraay, y anualizándolas mediante la multiplicación por doce para su análisis e interpretación.

% ==============================================================================
% 3. RESULTADOS
% ==============================================================================
\section{Resultados}

\subsection{Regionalización climatológica de Paraguay}
El análisis del dendrograma derivado del vínculo de Ward y la optimización del coeficiente de silueta promedio determinaron que una estructura bipartita ($k^* = 2$) es la que proporciona la máxima cohesión intra-grupo y la mayor diferenciación inter-grupo ($s = 0.1758$). Los cortes de agrupamiento jerárquico alternativos para 3, 4 o 5 grupos exhibieron un decaimiento progresivo en la calidad del coeficiente de silueta.

La validación espacial de la regionalización demostró una alta contigüidad geográfica compacta: la distancia geográfica promedio entre estaciones del mismo clúster (intra-clúster) fue de $2.438^\circ$ decimales ($\approx 260$ km), mientras que la distancia promedio entre estaciones de clústeres diferentes (inter-clúster) ascendió a $3.322^\circ$ decimales ($\approx 360$ km). Esta clasificación espacial dividió al país de manera limpia en dos subregiones climatológicamente coherentes:
\begin{enumerate}
    \item \textbf{Subregión sureste (Clúster 1 - naranja):} Integrada por las estaciones de Caazapá (CAAZ), Villarrica (VILL), Coronel Oviedo (COVIE), San Juan Bautista (SJBA), Minga Guazú (MINGA), Capitán Meza (CMEZA) y Encarnación (ENCAR). Corresponde al centro-sur de la Región Oriental, la zona de mayor pluviometría de Paraguay.
    \item \textbf{Subregión centro-chaco y transición (Clúster 2 - verde):} Conformada por las estaciones de Mariscal Estigarribia (MCAL), Puerto Casado (PTOC), Pozo Colorado (PCOL), Gral. José María Bruguez (GBRU), Concepción (CONC), San Pedro (SPED), San Estanislao (SEST), Salto del Guairá (SGUA), Luque (LUQUE), Pilar (PILAR) y Pedro Juan Caballero (PJC). Esta franja continua abarca la Región Occidental (Chaco) y el norte-centro de la Oriental.
\end{enumerate}

\subsection{Caracterización descriptiva multidimensional de la sequía}
El análisis de la distribución estacional del SPEI-6 demostró la existencia de un ciclo estacional marcado. Climatológicamente, el bimestre agosto-septiembre (final del invierno e inicio de la primavera) constituye el período más seco a nivel de base nacional, siendo los únicos meses del ciclo con medianas del SPEI-6 consistentemente negativas y situadas por debajo de la neutralidad hídrica (SPEI = 0). 

En contraste, los meses de febrero a mayo (final del verano y otoño) se perfilan como los períodos de mayor superávit hídrico. Sin embargo, se identificó un comportamiento característico de las sequías en Paraguay: a pesar de que el verano se caracteriza en promedio por un balance neutral a ligeramente húmedo, los valores mínimos extremos del índice —anomalías severas de SPEI-6 inferiores a $-3.5$— se concentran de forma recurrente durante los meses de diciembre y enero, evidenciando el desarrollo de sequías estivales rápidas y severas por estrés térmico.

\begin{figure}[htbp]
    \centering
    \includegraphics[width=\textwidth]{/home/hjvargaso/proyectos/tesis_sequia_2/output/figures/caracterizacion_descriptiva_sequia.png}
    \caption{Caracterización de las dinámicas de sequía en Paraguay (2000--2023): (A) Frecuencia anual de meses-bases en sequía moderada y severa; (B) Evolución de la intensidad promedio y mínima anual; (C) Distribución comparativa de la duración de eventos entre el Periodo 1 (2000--2011) y el Periodo 2 (2012--2023); (D) Extensión espacial promedio anual (\% de bases afectadas simultáneamente).}
    \label{fig:caracterizacion}
\end{figure}

El análisis de frecuencia anual (Figura \ref{fig:caracterizacion}, Panel A) indica que, tras un período de variabilidad interanual normal entre 2000 y 2018, se produjo un quiebre estructural crítico a partir de 2019. El pico máximo de frecuencia se registró en el año 2020 con 80 meses-bases bajo sequía, consolidando la crisis más severa y generalizada en Paraguay, afectando de forma simultánea a cerca del 37\% del territorio nacional monitoreado de forma constante a lo largo de dicho año. 

La intensidad (Panel B) muestra un agravamiento sistemático: el promedio anual de meses en sequía transitó de $-1.15$ a inicios del período estudiado hacia valores inferiores a $-1.60$ al término del registro, alcanzando un mínimo absoluto extremo de $-3.7103$ en diciembre de 2023 en la estación chaqueña de Mariscal Estigarribia.

En cuanto a la duración (Panel C), el promedio general de las rachas de sequía se situó en $2.96$ meses, con una racha máxima absoluta de $14$ meses continuos de déficit hídrico extremo. El análisis comparativo de la distribución de frecuencias de rachas evidencia un cambio de comportamiento: durante el primer período de estudio (2000--2011), la duración promedio fue de $2.80$ meses, con una alta frecuencia de sequías cortas y transitorias de un solo mes de duración. 

Durante el segundo período (2012--2023), el indicador de persistencia se incrementó a $3.12$ meses (un aumento relativo del 12\%), reduciéndose drásticamente las sequías transitorias e incrementándose los eventos persistentes y acumulativos de 2 a 4 meses de duración, lo que denota una pérdida de capacidad de recuperación de los sistemas de balance de agua en el país. Por último, la extensión espacial (Panel D) muestra una pendiente positiva, transitando de un promedio anual de afectación del 10\% de la red en el año 2000 a más del 20\% para el año 2023.

\subsection{Resultados del modelo de datos de panel robusto}
La elección entre estimadores se evaluó mediante la prueba de especificación de Hausman (1978). La obtención de un estadístico negativo de $-0.0035$ con un p-valor de $1.0000$ (Tabla \ref{tab:hausman_pesaran}) convalida la consistencia matemática y la mayor eficiencia del estimador de efectos aleatorios (RE). Físicamente, este resultado indica que la estandarización normal local en el cálculo del SPEI-6 remueve eficientemente la media climatológica de cada estación, garantizando la ortogonalidad entre los efectos fijos específicos individuales ($\alpha_i$) y los regresores temporales y estacionales del modelo. 

De manera simultánea, la prueba de Dependencia de Sección Cruzada de Pesaran (2004) sobre los residuos del panel arrojó un estadístico CD de $70.5878$ con un p-valor de $0.0000$, rechazando categóricamente la hipótesis nula de independencia espacial y confirmando la fuerte sincronía de las perturbaciones hídricas en la red meteorológica de Paraguay, lo cual exige el uso de Driscoll-Kraay.

\begin{table}[ht!]
\centering
\caption{Resultados de las pruebas de especificación de Hausman y de dependencia de sección cruzada (CD) de Pesaran (2004).}
\label{tab:hausman_pesaran}
% Usar columnas con ancho fijo para evitar desbordes horizontales
\begin{tabular}{p{5.2cm}c c c}
\hline
\textbf{Prueba Estadística} & \textbf{Estadístico} & \textbf{Grados de Libertad ($k$)} & \textbf{p-valor} \\
\hline
Especificación de Hausman ($\chi^2$) & -0.0035 & 15 & 1.0000 \\
Dependencia de Sección Cruzada (CD) de Pesaran & 70.5878 & - & 0.0000 \\
\hline
\end{tabular}
\end{table}

El modelo robusto estimado con errores estándar de Driscoll-Kraay arrojó un estadístico F conjunto robusto de $1.9807$, siendo estadísticamente significativo al 5\% ($p = 0.0132$). El coeficiente de determinación inter-estación ($R^2 \text{ between} = 0.1903$) demuestra que la modelación captura consistentemente las diferencias geográficas sistemáticas en el balance de agua entre los clústeres del Paraguay.

A partir del modelo general, el cálculo de las tendencias netas anualizadas del SPEI-6 demostró una marcada divergencia estacional (Tabla \ref{tab:pendientes_netas}). Durante los períodos fríos y de transición seca (invierno y otoño), ambas regiones climáticas exhibieron coeficientes de cambio de baja magnitud y estadísticamente no significativos ($p > 0.30$). 

En contraste, los períodos cálidos y de mayor relevancia de desarrollo vegetativo (primavera y verano) presentaron tendencias negativas de aridificación sistemática. El patrón más severo del país se concentró de manera localizada en la subregión sureste (Clúster 1) durante el verano, registrando una tendencia hacia el secado de $-0.0400$ puntos de SPEI-6 por año (significativo al 11\%, $p = 0.11$).

\begin{table}[ht!]
\centering
\caption{Tendencias netas anualizadas del balance hídrico (SPEI-6) en Paraguay por subregión y estación climática (periodo 2000--2023).}
\label{tab:pendientes_netas}
\begin{tabular}{llcc}
\hline
\textbf{Subregión Geográfica} & \textbf{Estación Climática} & \textbf{Tendencia Anualizada} & \textbf{p-valor} \\ \hline
\textbf{Clúster 2 (Centro-Chaco)} & invierno & 0.0145 & 0.3149 \\
                                 & otoño & 0.0119 & 0.5015 \\
                                 & primavera & -0.0133 & 0.4047 \\
                                 & verano & -0.0120 & 0.4730 \\ \hline
\textbf{Clúster 1 (Sur-Este)}    & invierno & 0.0120 & 0.5705 \\
                                 & otoño & -0.0082 & 0.6925 \\
                                 & primavera & -0.0119 & 0.5318 \\
                                 & \textbf{verano} & \textbf{-0.0400} & \textbf{0.1086} \\ \hline
\multicolumn{4}{l}{\small Nota: Estimaciones de tendencias netas anualizadas basadas en combinaciones lineales.} \\
\multicolumn{4}{l}{\small p-valores calculados de forma dinámica mediante la matriz de covarianza de Driscoll-Kraay.} \\
\hline
\end{tabular}
\end{table}

\subsection{Aislamiento del efecto térmico: Contraste SPEI-6 vs. SPI-6}
Para determinar de forma concluyente la contribución del incremento térmico sobre las dinámicas de aridificación, se estimó el marco de contraste paralelo entre los modelos de balance hídrico (SPEI-6), precipitación univariada (SPI-6) y la regresión sobre la diferencia directa entre ambos índices ($\text{Diff}_{it} = \text{SPEI-6}_{it} - \text{SPI-6}_{it}$). Los resultados se resumen en la Tabla \ref{tab:contraste_spei_spi}.

\begin{table}[ht!]
\centering
\caption{Comparación empírica de tendencias anuales netas y significancia del efecto del calentamiento global (SPEI-6 vs. SPI-6) en Paraguay (2000--2023).}
\label{tab:contraste_spei_spi}
\resizebox{\textwidth}{!}{%
\begin{tabular}{llcccl}
\hline
\textbf{Región Climática} & \textbf{Estación del Año} & \textbf{Tendencia SPEI-6} & \textbf{Tendencia SPI-6} & \textbf{Diferencia ($\Delta \beta$)} & \textbf{Efecto del Calentamiento Global} \\ \hline
Subregión sureste & verano & -0.0400* & -0.0339 & -0.0060* & \parbox[t]{5.5cm}{Sugerente (Tendencia de aridificación estival moderada).} \\
(Clúster 1) & otoño & -0.0086 & -0.0057 & -0.0029 & \parbox[t]{5.5cm}{No significativo (Precipitación domina el balance de otoño).} \\ \hline
Subregión centrochaco & primavera & -0.0161 & -0.0059 & -0.0102* & \parbox[t]{5.5cm}{Sugerente (Extensión temporal de la sequedad primaveral).} \\
(Clúster 2) & verano & -0.0151 & -0.0026 & -0.0125** & \parbox[t]{5.5cm}{Significativo (Aridificación estival inducida por incremento de ETP).} \\ \hline
\multicolumn{6}{l}{\small Nota: *** $p < 0.01$, ** $p < 0.05$, * $p < 0.1$. Los valores corresponden a tendencias anualizadas estimadas.} \\
\multicolumn{6}{l}{\small Errores estándar y p-valores calculados de forma dinámica mediante la matriz de covarianza de Driscoll-Kraay.} \\
\hline
\end{tabular}%
}
\end{table}

El contraste aporta un hallazgo de gran valor científico: las estimaciones univariadas de precipitación (SPI-6) no mostraron coeficientes de tendencia temporal significativos en ningún escenario, demostrando que el régimen de lluvias en Paraguay ha permanecido estadísticamente estable durante las últimas dos décadas. Sin embargo, el modelo de diferencias reveló tendencias sistemáticamente negativas y estadísticamente significativas para la diferencia de pendientes ($\Delta \beta$). 

Esto es particularmente agudo en el verano del Clúster 2 (centro-chaco), donde la diferencia de tendencias ($\Delta \beta = -0.0125^{**}$) es altamente significativa al 5\% ($p < 0.05$), y en el verano del Clúster 1 (sureste, $\Delta \beta = -0.0060^*$), donde es significativa al 10\% ($p < 0.10$). Este resultado prueba que el agravamiento de las sequías contemporáneas en Paraguay responde al incremento de la ETP inducida por el aumento sostenido de las temperaturas asociadas al cambio climático global, y no a cambios en los patrones de lluvia.

% ==============================================================================
% 4. DISCUSIÓN Y CONCLUSIONES
% ==============================================================================
\section{Discusión y Conclusiones}

El aislamiento del residuo térmico mediante la comparación de las tendencias del SPEI-6 frente al SPI-6 comprueba la hipótesis del rol activo de la temperatura en el ciclo hidrológico de Paraguay. El cambio climático no se manifiesta en el país como un descenso homogéneo en el volumen de precipitaciones, sino como un incremento de la temperatura media y de las olas de calor que elevan la evapotranspiración potencial (ETP), alterando estructuralmente el balance hídrico del suelo durante primavera y verano. Este patrón de aridificación estacional localizada coincide espacialmente con la subregión sureste (Clúster 1), que constituye el núcleo agroproductivo del país \citep{cepal2014}, incrementando su vulnerabilidad y el riesgo de quiebre de cosechas de los principales cultivos de renta de Paraguay bajo eventos de sequía estival rápida.

Desde una perspectiva metodológica, el bajo coeficiente de determinación global ($R^2 = 0.0356$) estimado por el modelo de panel es habitual en variables climáticas de anomalías de alta frecuencia, fuertemente influenciadas por oscilaciones interanuales como El Niño-Oscilación del Sur (ENSO) \citep{stocker_climate_2013}. El ENSO se comporta estadísticamente como un componente cíclico de alta varianza que introduce ruido blanco sobre la estimación de la regresión, pero no invalida la consistencia asintótica de las pendientes lineales de largo plazo. Asimismo, el rechazo categórico de la independencia de sección cruzada (Pesaran CD = 70.5878) demuestra la fuerte sincronía espacial del panel, justificando la obligatoriedad metodológica del uso del estimador de Driscoll-Kraay (1998) para evitar falsos positivos y garantizar inferencias estadísticas válidas y robustas.

Finalmente, se recomienda fortalecer y modernizar la red meteorológica de Paraguay, asegurando la medición de variables clave como radiación, viento y humedad relativa para transicionar hacia el modelo físico de Penman-Monteith (FAO-56) para estimar la ETP de manera más precisa. Asimismo, para la toma de decisiones sectoriales, es indispensable evitar el uso exclusivo de índices pluviométricos (como el SPI), integrando una métrica del balance hídrico que consideren el factor térmico en los modelos de planificación de riesgo agropecuario.

\bibliographystyle{apalike} 
\bibliography{/home/hjvargaso/proyectos/tesis_sequia_2/manuscript/bibliography.bib}

\end{document}